\scanpage{110}
% archive sheet number: 109
% source 01-04, top
Si une réalisation est acceptable, \deleted{\unsure{les}} \unsure{prenons|prenant}
quatre pts distincts $s_0,s_1,s_2,s_3$ dont les trois premiers sont sur
une $=$ face, leurs réal. géom. sont aff\textsuperscript{t}
\unsure{indépendantes|indépendants}.
% source 05-08, upper middle, revised definition
\deleted{Considérons} Une représentation \inserted{est dite}
\deleted{\uuline{acceptable}} \deleted{\unsure{dans le}}
\inserted{\uline{régulière}} si \inserted{elle est acceptable et si}
pour \unsure{tt|tout} $g\in\Gamma$ $\exists!\,g_E$ qui est un affine de $E$
qui \unsure{le|la} prolonge.
% source 09-11, middle
\deleted{\unsure{Pour en faire}}\inserted{C'est une}
\deleted{\unsure{on définit} les repr. acceptables}
propriété de l'application $\unsure{\rho_S|\rho_s}$ en fait, et
n'implique pas $\rho_S(S)$ engendre affinement $E$.
% source 12, left slanted margin linked to the definition
% slanted
\margin{$\times$ \uline{NB} la condition ``acceptable'' revient à exiger que
$\rho_S$ \unsure{soit} un \unsure{point|plgt}
$\unsure{\operatorname{Conv}S|\operatorname{Conf}S}\to E$, ou encore que
$\rho_S(S)$ ne soit pas contenu dans une droite...}
% source 13, middle
\uline{Prop.} Une représentation régulière est \uline{fidèle}.
% source 14-16, lower middle
Soit $\deleted{(\unsure{\delta_0|\sigma_0},\delta_1,\delta_2)}$
$(\sigma_0,\sigma_1,\sigma_2)\in\Gamma$ les symétries fond. associées
au repère, $\deleted{\unsure{\varphi|s}}s_0$ le sommet origine,
% source 17, lower middle
\[
s_1=\sigma_0s_0,\quad s_2=\sigma_1s_1,\quad s_3=\sigma_2s_2
\]
% source 18-20, lower middle
\deleted{Je dis que}
Je dis que $s_0,s_1,s_2,s_3$ sont aff\textsuperscript{t} indép. :
On sait déjà par hyp. que $s_0,s_1,s_2$ le sont.
% source 21-24, bottom
Si on avait $s_3\in P:=\operatorname{Pl}(s_0,s_1,s_2)$, il est visible que
$P$ --- qui est déjà invariant par $\sigma_0,\sigma_1$ --- serait invariant
par $\sigma_2$ $\deleted{(\unsure{\sigma_2|\sigma_1})}$ puisque
\[
\sigma_2(P)=\operatorname{Pl}(\sigma_2(s_0)=s_0,\ \sigma_2(s_1)=s_1,\
\sigma_2(s_2)=s_3)=P
\]
% source 25, bottom, continuation stops at the end of this scan
Ceci étant, on peut \deleted{\unsure{faire}} choisir $s_0,s_1,s_2,s_3$
